\section{Discussion}\label{sec:discussion}

\paragraph{What the evidence supports}
Three findings are supported by paired statistics over 16 seeds. First, when the budget is binding and the data are strongly label-skewed across clients (PAMAP2, $\phi\le0.25$), letting clients choose how much of the network to train, under energy prices and a coverage constraint, gives a much better final model than FedAvg-type training with random or utility-based selection (+6 to +27 pp). Second, when the labels are similar across clients (UCI HAR) the same control gives the same accuracy within 1--3 pp for $1.6$--$2.2$ times less energy and with far fewer drained devices. Third, the pieces that matter are the energy queue, depth control and epoch control; the coverage queue matters for the deep exits, and precision control matters little unless radio energy is expensive relative to computation (Section~\ref{sec:sens}). The theory is consistent with the coverage experiment, although the relation is noisy and the bound is loose.

\paragraph{What it does not support}
EcoFed does not dominate. Static-Depth---a one-line rule that gives each device the deepest depth that fits its budget---reaches the best checkpoint accuracy on PAMAP2 at the tightest budget ($+7.8$ pp over EcoFed) and ties EcoFed at the loosest, and its energy to a target accuracy is statistically indistinguishable from EcoFed's. What EcoFed adds over that rule is accuracy at the \emph{end} of the run, usable deep exits, protection of the batteries and a single knob, $\lambda$, that trades energy for accuracy (Section~\ref{sec:pareto}). On UCI HAR the most accurate method is single-exit FedAvg, which beats EcoFed by 2--3 pp at the best checkpoint; a multi-exit model is only worth its (small) extra training cost if the deployment-time savings of Section~\ref{sec:lifecycle} are wanted.

\paragraph{Why depth control may help when labels are skewed}
A plausible explanation, which we did not test in isolation, is breadth of participation: cheap shallow updates let many clients---and therefore many activities---contribute in every phase of the run, while full-depth updates exhaust a few clients and then let the fleet shrink to those that can afford them. The accuracy collapse of the energy-blind baselines after their peak on PAMAP2 and the near-parity on UCI HAR, where every client holds every label, are both consistent with this reading.

\paragraph{Implications for sustainable edge AI}
The lifecycle analysis reframes where the energy goes. In our model, a fleet that serves a wearable classifier continuously spends as much energy on inference as it spent on training within one to two days, and for the cascades and the PAMAP2 and UCI HAR models studied here 18 to 270 times as much over a month (Table~\ref{tab:life}). Two consequences follow. A federated training controller should be judged not only by the Joules it saves during training---which in absolute terms are small---but by whether it produces models that can be served cheaply: EcoFed's shallow exits are confident enough that 55--98\% of windows leave at the first block, while the exits of the model trained with a plain multi-exit FedAvg are not. And the training cost of making a model multi-exit is repaid within minutes of operation, so there is no energy argument for deploying the single-exit network, even where single-exit training reaches a slightly higher accuracy.


\paragraph{Practical guidance}
(i) Estimate the budget regime first: if the energy per participation of the full model is within a factor of a few of the per-round budget, depth and epoch control pay off; if energy is abundant, they only save energy. (ii) If the data are close to IID across clients, expect little accuracy to be lost by shallow, cheap updates, and use the savings. (iii) If the operator can stop training at a validation-selected checkpoint and does not need deep exits, a static capability-based depth is a strong, simple baseline. (iv) Keep the energy accounting closed-loop: EcoFed's queues use energy reported by the devices, which is what makes it robust to an incorrect energy model.

\section{Limitations and threats to validity}\label{sec:limits}
\begin{itemize}[leftmargin=1.4em,itemsep=1pt]
\item \textbf{Energy is modelled, not measured.} The device constants are assumptions and the paper contains no power measurement. Section~\ref{sec:sens} shows that conclusions are robust to radio cost, to multiplicative errors of up to 50\% (log-normal $\sigma=0.5$) in the controller's belief, and to an unmodelled per-sample overhead fitted from CPU latency; but a power-metered Raspberry Pi/Jetson study is the necessary next step (the measurement and fitting harness is provided, \ref{app:hardware}, but was not run), and the absolute Joules reported here should not be quoted as measurements.
\item \textbf{Scale.} The fleets have 16--20 clients and the network 27\,784 parameters; the selection is 5 clients per round over 80 rounds. We do not know how the ranking of methods changes with thousands of clients, deeper networks, or larger models; the controller overhead measurements (Section~\ref{sec:overhead}) say only that the algorithm itself scales.
\item \textbf{Two datasets from one domain.} Both are wearable inertial benchmarks. Results on audio or vision, and on datasets with feature rather than label skew, may differ.
\item \textbf{Tuning asymmetry.} Every method was tuned on separate seeds and validation subjects with a two-stage search, but EcoFed has more knobs (14 trials) than most baselines (3--7). Several selected values lie on the edge of their grid (for example the learning rate, and the depth exponent $\kappa_d=0.5$ for EcoFed), so all methods may be slightly under-tuned. The setting $\phi=1$ reuses the $\phi=0.25$ configuration. A pilot on tuning seeds showed that randomised instead of top-$m$ selection and $\kappa_d=0.25$ changed the validation AUC by between $-2.4$ and $+2.5$ points depending on the setting; we kept the pre-specified configuration rather than re-tune after seeing test results.
\item \textbf{Theory versus practice.} Theorem~\ref{thm:main} assumes iterate-independent participation, bounded gradients and an unbiased Horvitz--Thompson aggregator; the experiments use utility-based selection and sample-normalised block averages. Proposition~\ref{prop:lyap} is stated for an idealised system. The coverage experiment supports the qualitative prediction only.
\item \textbf{Exit weights.} We weight exits proportionally to $j^2$ for all methods. This favours the last exit and handicaps the shallow exits of full-depth methods, which affects the cascade comparison (Table~\ref{tab:exits}); we did not test other weightings beyond the pilot of \ref{app:extra}.
\item \textbf{Baselines.} Static-Depth and Oort-style are our simplified implementations of ideas from the literature, not reproductions of DepthFL/HeteroFL/ScaleFL or the Oort system.
\item \textbf{Validation-selected checkpoint.} The best-checkpoint metric assumes a small validation set at the server; the final-accuracy metric does not.
\item \textbf{Not modelled.} Stragglers and dropouts, wireless channel dynamics, charging, privacy mechanisms and security, and the energy of the server.
\end{itemize}

\section{Conclusion}\label{sec:conclusion}
We studied federated learning of multi-exit networks on battery-limited fleets, where energy is spent twice---when training and when serving. A convergence bound shows that letting weak devices train only shallow blocks creates an error floor controlled by how often each block is trained; EcoFed turns this into a coverage constraint and combines it with per-device energy queues and a shadow price in a controller that costs $O(N|\mathcal C|)$ per round. On two wearable benchmarks it improves the final accuracy of the federated model substantially when labels are strongly skewed and the budget binds, matches the accuracy of simpler methods at much lower energy when labels are similar, protects device batteries, and degrades gracefully when the energy model is wrong or radio energy dominates. It does not dominate a simple capability-based depth rule at the tightest budget, and single-exit training remains more accurate on the easy dataset. Future work should (i) calibrate the energy model on metered devices and repeat the study on hardware, (ii) scale to larger fleets and models, including audio and language workloads, (iii) learn the utility of depth online instead of tuning its exponent, and (iv) couple the training-time controller with the deployment-time exit policy so that the energy of both phases is optimised together.
